\documentclass[11pt,a4paper]{article}

\usepackage[T1]{fontenc}
\usepackage[utf8]{inputenc}
\usepackage[lithuanian]{babel}
\usepackage{newtxtext,newtxmath}

\usepackage[
  landscape,
  left=2cm,
  right=1.2cm,
  top=1.25cm,
  bottom=1cm
]{geometry}

\usepackage{enumitem}
\usepackage{xcolor}
\usepackage{tabularray}
\usepackage{ragged2e}
\usepackage{microtype}

\pagestyle{empty}

\setlength{\parindent}{0pt}
\setlength{\parskip}{0pt}

\setlist[itemize]{
  leftmargin=1.45em,
  label=\textbullet,
  itemsep=0pt,
  topsep=0pt,
  parsep=0pt,
  partopsep=0pt
}

\setlist[enumerate]{
  leftmargin=1.3em,
  label=\arabic*.,
  itemsep=0pt,
  topsep=0pt,
  parsep=0pt,
  partopsep=0pt
}

\definecolor{tablegray}{gray}{0.82}

\newcommand{\DocHeading}[1]{%
  \vspace{0.75\baselineskip}%
  {\large #1}\par
  \vspace{0.15\baselineskip}%
}

\begin{document}

\sloppy

% =========================================================
% 1 PUSLAPIS
% =========================================================

\begin{flushright}
\begin{minipage}{8.5cm}
PATVIRTINTA\\
Vilniaus licėjaus direktoriaus\\
2023 m. sausio 31 d. įsakymu V1-21
\end{minipage}
\end{flushright}

\vspace{0.55cm}

\begin{center}
{\large VILNIAUS LICĖJUS}\par

\vspace{0.25cm}

{\large 2025--2026 m. m. MATEMATIKOS MODULIO PROGRAMA}
\end{center}

\vspace{0.45cm}

\begin{tblr}{
  colspec={
    Q[l,wd=5.4cm]
    Q[l,wd=8.5cm]
  },
  rowsep=0pt,
  colsep=0pt
}
DALYKAS & MATEMATIKA \\
KLASĖ & I \\
PAMOKŲ SKAIČIUS & 36 (1 pamoka per savaitę) \\
\end{tblr}

\vspace{0.38cm}

\textbf{TIKSLAS} -- sudaryti galimybę kiekvienam mokiniui, mokantis
matematikos, ugdytis matematinį ir statistinį raštingumą, kuris šiame
dokumente suprantamas kaip įgytas gebėjimas matematiškai samprotauti ir
taikyti įgytas kompetencijas, sprendžiant įvairias realias, aktualias ir
mokiniams suprantamas problemas.

\DocHeading{UŽDAVINIAI}

Siekdami tikslo mokiniai:

\begin{itemize}

  \item tinkamai ir tikslingai vartoja matematinius faktus; sklandžiai
  atlieka matematines procedūras; įgytas žinias sieja tarpusavyje,
  sistemina, struktūruoja; įžvelgia matematikos ryšius su kitais
  dalykais;

  \item įvairiuose kontekstuose taiko indukcinį ir dedukcinį,
  kiekybinį ir statistinį samprotavimą; remiasi žiniomis, logika ir
  patikimais argumentais, formuluodami, analizuodami, įrodinėdami
  teiginius, spręsdami uždavinius, darydami išvadas ar vertinimus;

  \item bendradarbiaudami su kitais, nagrinėja įvairiomis formomis
  pateiktus matematinius pranešimus, dalyvauja diskusijose apie
  komunikavimo tikslą, adresatą, pranešimu perteikiamų minčių
  tikslumą, logiškumą, pagrįstumą, išsamumą, glaustumą;

  \item yra nusiteikę ir įdeda pastangų matematikos mokymosi kliūtims
  įveikti; tikslingai planuoja ir organizuoja mokymosi veiklą;
  siekdami mokytis matematikos ir ją pažinti, turi žinių, gebėjimų ir
  polinkį naudotis skaitmeninėmis technologijomis;

  \item įgytas matematines kompetencijas ir supratimą apie bendrą
  problemų sprendimo procesą kūrybiškai pritaiko įvairiuose realiuose,
  aktualiuose ir mokiniams suprantamuose kontekstuose; reflektuoja
  savo žinias, gebėjimus, samprotavimo veiklą ir jos rezultatus.

\end{itemize}

\DocHeading{MOKYMO IR MOKYMOSI PRIEMONĖS}

\begin{enumerate}

  \item Matematika 9 klasei. 1, 2 vadovėlio dalys
  (serija „Horizontai“). Jūratė Gedminienė, Irena Šukienė,
  Daiva Riukienė, Jolanta Jačiauskaitė, Ingrida Brazauskienė.
  Kaunas: „Šviesa“, 2023. (prieiga „Eduka klasėje“).

  \item Matematika. Vadovėlis 9 klasei. 1, 2 vadovėlio dalys
  (serija „TEMPUS“). Jolanta Knyvienė, Antanas Apynis,
  Ingrida Brazauskienė, Ramunė Dranseikienė, Nomeda Zdanienė.
  Kaunas: „Šviesa“, 2020.

  \item Matematika. Vadovėlis gimnazijų I klasei.
  1, 2 vadovėlio dalys. Alvyda Ambraškienė, Rita Belkevičienė,
  Rita Grigelienė, Birutė Vosylienė, Milda Vosylienė.
  Kaunas: „Šviesa“, 2008.

\end{enumerate}

\DocHeading{VERTINIMAS}

\newpage


% =========================================================
% 2 PUSLAPIS
% =========================================================

\setlength{\parindent}{0.9cm}

Mokinių pasiekimai vertinami vadovaujantis Vilniaus licėjaus mokinių
pasiekimų ir pažangos vertinimo tvarkos aprašu, patvirtintu direktoriaus
2024-06-27 įsakymu Nr. V1-112. Atsižvelgiant į pamokos uždavinius,
nuolat taikomas formuojamasis vertinimas.

Mokinių pasiekimai matematikos pamokose vertinami tarpinėmis įskaitomis,
kurių per pusmetį reikia surinkti dvi. Pirmoji tarpinė įskaita rašoma po
pirmųjų devynių pamokų, jeigu mokinys atliko bent penkių pamokų klasės
darbą, aktyviai dalyvavo pamokose. Antroji tarpinė įskaita paskutinę
pusmečio pamoką, jeigu mokinys per paskutines devynias pamokas atliko
bent penkių pamokų klasės darbą, aktyviai dalyvavo pamokose.

\setlength{\parindent}{0pt}

\vspace{0.28cm}
\newpage
{\large MOKYMO IR MOKYMOSI TURINYS}\par 

\vspace{0.5cm}

\begingroup

\fontsize{9.5}{10.1}\selectfont

\SetTblrInner{
  rowsep=0.9pt,
  colsep=4pt
}

\begin{tblr}{
  width=\textwidth,
  colspec={
    |Q[l,wd=10.15cm]
    |Q[c,wd=2.15cm]
    |Q[l,wd=6.55cm]
    |X[l]|
  },
  hlines,
  vlines,
  cells={valign=t}
}

% =========================================================
% 1 TEMA
% =========================================================

\SetCell[c=3]{l,bg=tablegray,font=\bfseries}
\textbf{1 Tema. Laipsniai ir šaknys}
&
&
&
\SetCell{l,bg=tablegray,font=\bfseries}
Valandų skaičius: 4
\\

\SetCell{c,font=\bfseries}
Pamokos tema
&
\SetCell{c,font=\bfseries}
Valandos
&
\SetCell{c,font=\bfseries}
Kompetencijos ir veiklos
&
\SetCell{c}
\shortstack{
\bfseries Pastabos\\
(vertinimas, edukacinės išvykos,\\
papildomos priemonės ir kt.)
}
\\

1. Laipsnių savybės ir reiškinių su laipsniais pertvarkymas
&
2
&
\SetCell[r=3]{l}
Pažinimo kompetencija
(sisteminamos žinios apie laipsnių ir šaknų savybes, pasirenkami
tinkami reiškinių pertvarkymo būdai, atliekami skaitinių ir raidinių
reiškinių skaičiavimai).

Skaitmeninė kompetencija
(skaičiuotuvas tikslingai naudojamas skaičiavimų rezultatams
patikrinti ir labai dideliems bei labai mažiems skaičiams nagrinėti).
&
\SetCell[r=3]{l}
\\

2. Kvadratinė ir kubinė šaknis. Šaknų savybės
&
1
&
&
\\

3. Standartinė skaičiaus išraiška ir jos taikymas
&
1
&
&
\\


% =========================================================
% 2 TEMA
% =========================================================

\SetCell[c=3]{l,bg=tablegray,font=\bfseries}
\textbf{2 Tema. Lygtys ir nelygybės}
&
&
&
\SetCell{l,bg=tablegray,font=\bfseries}
Valandų skaičius: 3
\\

1. Tiesinės lygtys ir jų sprendinių analizė
&
2
&
\SetCell[r=2]{l}
Pažinimo kompetencija
(kartojami tiesinių lygčių ir nelygybių sprendimo algoritmai,
nagrinėjami atvejai, kai sprendinių nėra arba jų yra be galo daug).

Komunikavimo kompetencija
(nuosekliai užrašomi sprendimo žingsniai, paaiškinami atliekami
pertvarkymai ir taisyklingai užrašoma sprendinių aibė).
&
\SetCell[r=2]{l}
\\

2. Tiesinės nelygybės ir jų sprendinių aibės
&
1
&
&
\\


% =========================================================
% 3 TEMA
% =========================================================

\SetCell[c=3]{l,bg=tablegray,font=\bfseries}
\textbf{3 Tema. Tekstinių uždavinių kartojimas}
&
&
&
\SetCell{l,bg=tablegray,font=\bfseries}
Valandų skaičius: 8
\\

1. Uždaviniai, sprendžiami sudarant lygtį
&
2
&
\SetCell[r=5]{l}
Pažinimo kompetencija
(analizuojama uždavinio sąlyga, išskiriami žinomi ir nežinomi
dydžiai, sudaromas matematinis situacijos modelis ir įvertinamas
gauto atsakymo prasmingumas).

Kūrybiškumo kompetencija
(lyginami skirtingi uždavinio sprendimo būdai, pasirenkama tinkama
strategija ir ieškoma efektyvesnio sprendimo).

Pilietiškumo kompetencija
(nagrinėjamos kasdienio gyvenimo, finansinio raštingumo, judėjimo
ir darbo situacijos).
&
\SetCell[r=5]{l}
\\

2. Proporcijų ir procentų uždaviniai
&
2
&
&
\\

3. Judėjimo ir kelio uždaviniai
&
2
&
&
\\

4. Bendro darbo ir našumo uždaviniai
&
1
&
&
\\

5. Mišrūs realaus konteksto tekstiniai uždaviniai
&
1
&
&
\\


% =========================================================
% 4 TEMA
% =========================================================

\SetCell[c=3]{l,bg=tablegray,font=\bfseries}
\textbf{4 Tema. Planimetrijos kartojimas}
&
&
&
\SetCell{l,bg=tablegray,font=\bfseries}
Valandų skaičius: 12
\\

1. Kampai tarp susikertančių ir lygiagrečių tiesių
&
1
&
\SetCell[r=7]{l}
Pažinimo kompetencija
(kartojamos pagrindinės plokštumos geometrijos sąvokos ir teoremos,
atpažįstamos geometrinių figūrų savybės bei ryšiai tarp jų).

Komunikavimo kompetencija
(geometriniai teiginiai pagrindžiami nuosekliais argumentais,
taisyklingai vartojamos geometrinės sąvokos ir žymėjimai).

Kūrybiškumo kompetencija
(sprendžiant uždavinius pasirenkamos tinkamos teoremos, atliekamos
papildomos konstrukcijos ir derinami keli geometriniai faktai).
&
\SetCell[r=7]{l}
\\

2. Trikampiai, jų rūšys ir pagrindinės savybės
&
2
&
&
\\

3. Pitagoro teorema ir statinio prieš $30^\circ$ kampą savybė
&
2
&
&
\\

4. Trikampio pusiaukampinė, pusiaukraštinė ir aukštinė
&
1
&
&
\\

5. Trikampių lygumo požymių taikymas
&
2
&
&
\\

6. Keturkampiai, jų rūšys ir pagrindinės savybės
&
2
&
&
\\

7. Lygiagretainis, trapecija ir jų savybių taikymas
&
2
&
&
\\

\end{tblr}

\endgroup


% =========================================================
% 3 PUSLAPIS
% =========================================================

\newpage

\begingroup

\fontsize{10.35}{11.1}\selectfont

\SetTblrInner{
  rowsep=1.5pt,
  colsep=4pt
}

\begin{tblr}{
  width=\textwidth,
  colspec={
    |Q[l,wd=10.15cm]
    |Q[c,wd=2.15cm]
    |Q[l,wd=6.55cm]
    |X[l]|
  },
  hlines,
  vlines,
  cells={valign=t}
}

% =========================================================
% 5 TEMA
% =========================================================

\SetCell[c=3]{l,bg=tablegray,font=\bfseries}
\textbf{5 Tema. Stereometrijos kartojimas}
&
&
&
\SetCell{l,bg=tablegray,font=\bfseries}
Valandų skaičius: 5
\\

1. Kubas ir stačiakampis gretasienis
&
1
&
\SetCell[r=3]{l}
Pažinimo kompetencija
(kartojamos pagrindinės erdvinių kūnų sąvokos, atpažįstami jų
elementai, nagrinėjamos briaunainių ir sukinių savybės).

Erdvinio mąstymo gebėjimai
(plokštumoje pateiktas brėžinys siejamas su erdviniu kūnu,
analizuojamos kūnų išklotinės ir jų elementų tarpusavio ryšiai).

Kūrybiškumo kompetencija
(žinomos formulės ir geometrinės savybės taikomos sprendžiant
erdvinės geometrijos ir praktinio konteksto uždavinius).
&
\SetCell[r=3]{l}
\\

2. Stačioji prizmė ir piramidė
&
2
&
&
\\

3. Kūgis ir ritinys
&
2
&
&
\\


% =========================================================
% REZERVINĖS VALANDOS
% =========================================================

\SetCell{l,bg=tablegray,font=\bfseries}
Rezervinės ir apibendrinimo valandos
&
\SetCell{c,bg=tablegray,font=\bfseries}
4
&
\SetCell[c=2]{l,bg=tablegray}
&
\\


% =========================================================
% IŠ VISO
% =========================================================

\SetCell{r,font=\bfseries}
Iš viso
&
\SetCell{c,font=\bfseries}
36
&
\SetCell[c=2]{l}
&
\\

\end{tblr}

\endgroup


% =========================================================
% PARAŠAI
% =========================================================

\vspace{0.7cm}

\begin{tblr}{
  width=\textwidth,
  colspec={X[l] X[r]},
  colsep=0pt,
  rowsep=0pt
}
Parengė matematikos mokytojas
&
Andrius Berniukevičius
\end{tblr}

\vspace{1.25cm}

APTARTA\\
Matematikos metodinėje grupėje\\
Protokolo Nr. 1\\
2026 -- 08 --\\[0.42cm]

\rule{3.7cm}{0.4pt}\\[-1pt]
\hspace{0.75cm}{\small (parašas)}\\

Metodinės grupės pirmininkė\\
Inga Vaičiulė

\vspace{0.5cm}

SUDERINTA\\[0.42cm]

\rule{3.7cm}{0.4pt}\\[-1pt]
\hspace{0.75cm}{\small (parašas)}

\end{document}